\begin{abstract}
In December 2021 the Securities and Exchange Board of India (SEBI) suspended
exchange-traded derivatives in seven agricultural commodities. We evaluate whether
the suspension delivered its stated price-stability objective using district-level
mandi spot prices from CEDA/Agmarknet, realized-volatility panels, difference-in-
differences, Synthetic DiD, and synthetic control. An inherited claim held that the
suspension raised spot volatility by roughly $8$--$10\%$. That claim is refuted.
The current food-donor rerun estimates an aggregate DiD effect of $-9.8\%$, but it
is not statistically significant under honest few-cluster inference (CR1 $p=0.145$;
wild bootstrap $p=0.153$). The lower-volatility signal is concentrated in chana:
Synthetic DiD gives $-38.7\%$ ($z=-2.05$), and commodity-level synthetic control gives
$-37.3\%$ at the placebo-resolution floor. Wheat collapses to approximately zero under
Synthetic DiD, consistent with MSP/export-ban confounding. The honest conclusion is
therefore not that the ban robustly reduced volatility, but that it did not raise spot
volatility and that any lower-volatility evidence is chana-centered and donor-sensitive.
The full food/oilseed donor pull and additional inference upgrades remain pending.
\end{abstract}

\section{Introduction}
\label{sec:intro}

\subsection{The policy and the puzzle}

On 19--20 December 2021, SEBI directed India's commodity exchanges to suspend new
derivatives contracts and fresh positions in wheat, paddy, soybean complex, crude
palm oil and moong, while earlier restrictions on chana and mustard continued. The
policy was justified as a price-stability measure during a food-inflation episode.
The natural empirical question is whether shutting futures markets calmed the
underlying spot markets or removed a useful price-discovery and hedging venue.

Theory does not settle the sign. Futures markets can stabilize spot prices by
aggregating information and coordinating storage, but they can also transmit noise,
speculation or financial shocks. India's episode is useful because several close
commodities continued trading and can serve as donors, while the suspension itself
was large, persistent and policy-salient.

\subsection{What this paper finds}

The key result is a calibrated rejection of the inherited positive-volatility claim.
On the current food-cereal donor panel, the aggregate DiD is $-9.8\%$ with CR1
$p=0.145$ and wild-bootstrap $p=0.153$. The sign is negative, not positive, but the
aggregate result is not statistically robust. Compared with the earlier clean-core
run ($-20.7\%$), adding food-cereal donors halves the estimate and removes aggregate
significance.

The lower-volatility evidence is concentrated in chana, the cleanest treated
commodity. Synthetic DiD estimates chana at $-38.7\%$ with $z=-2.05$, and synthetic
control estimates $-37.3\%$ with the placebo test at its resolution floor. Wheat no
longer supports the story: under food-donor Synthetic DiD it is $+0.4\%$, which is
consistent with MSP, export-ban and stock-limit confounding.

\subsection{Contribution}

The paper contributes in four ways. First, it stress-tests and rejects the inherited
claim that the suspension increased spot volatility. Second, it shows that donor choice
is first-order: a clean-core donor pool and a food-staple donor pool imply different
strengths of inference. Third, it documents a commodity-level result in which chana,
not the aggregate, carries the lower-volatility signal. Fourth, it keeps the benefit
arm and cost arm separate: lower or unchanged spot volatility is not a welfare verdict
if price discovery, hedging and market integration are impaired.

The rest of the paper is organized as follows. Section~\ref{sec:institutions} describes
the suspension and institutional context. Section~\ref{sec:data} explains the spot-price
panel and donor construction. Section~\ref{sec:methods} lays out the empirical strategy.
Sections~\ref{sec:results} and~\ref{sec:robustness} report the current estimates and
robustness. Section~\ref{sec:discussion} discusses interpretation, limitations and policy
implications.
